\documentclass[11pt,a4paper]{article}
\usepackage[margin=25mm]{geometry}
\usepackage{amsmath,amssymb,amsthm,hyperref}
\setlength{\emergencystretch}{2em}
\newtheorem{lemma}{Lemma}
\newtheorem{theorem}[lemma]{Theorem}
\newcommand{\gen}[1]{\langle#1\rangle}
\title{The residual centralizer at a singleton end}
\author{Human Compression project: L044 continuation}
\date{}
\begin{document}\maketitle
Let $G=\gen{\rho_0,\ldots,\rho_{r-1}}\leq A_n$ be a transitive
string C-group, $n\geq12$, with all maximal omissions intransitive.
Suppose $G_0$ has orbits $\{a\}$ and $B$, and write the unique
crossing as $ab$. Put $X=\Omega\setminus\{a,b\}$,
$K=G_{>1}$, and $\rho_0=(a\ b)\delta$.
Then $K$ fixes $a,b$, preserves $X$, and commutes with the odd
involution $\delta$ on $X$. Neither transitivity nor a full-support
assumption on $K$ is made.

\begin{lemma}[An actual orbit-deficit inequality]\label{deficit}
Let $t$ be the number of two-point $\delta$-orbits and $f$ its
number of fixed points in $X$. Let $z$ be the number of orbits
of $\gen{K,\delta}$ on $X$, including singleton orbits. Then
\[
                    r-2\leq t+f-z.                    \tag{1}
\]
Moreover the distinguished images of the $r-2$ generators of $K$
in the permutation action on the $t+f$ orbits of $\delta$
form an independent generating set.
\end{lemma}
\begin{proof}
Commutation makes the partition into $\delta$-orbits invariant
under $K$. For a distinguished generator $s=\rho_j$,
$j\geq2$, let $K^-_j$ be generated by the other distinguished
$K$ generators. Suppose its image has the same orbits on the
$\delta$-blocks as the full image of $K$.
Adjoining $\delta$ makes the point orbits of $K^-_j$ on $X$
exactly the unions of blocks in these block orbits: inside
every two-point block $\delta$ supplies both points.
Thus $\gen{K^-_j,\rho_0}$ and $\gen{K,\rho_0}$ have exactly
the same point orbits on $\Omega$, including the orbit $\{a,b\}$.
Adjoining the same permutation $\rho_1$ to orbit-equivalent
groups leaves them orbit-equivalent, since the resulting
partition is the equivalence closure of the old partition
and the edges of $\rho_1$. This would make $G_j$ transitive,
a contradiction. Therefore deleting any distinguished image
strictly refines its orbit partition. In particular no image
is generated by the others, proving independence.

Choose for each distinguished image one of its edges joining
different orbits after that image is deleted. The selected
edges form a forest on the $t+f$ blocks: a cycle would give
an alternate path avoiding the label of any one of its edges,
contradicting that edge's eligibility. The forest has at least
$z$ components because every selected edge stays inside a
full block-action orbit. It has $r-2$ edges, giving (1).
\end{proof}

\begin{theorem}[No fixed-point-free residual exception]\label{nofixed}
If $\delta$ has no fixed point on $X$, then
$r\leq\lfloor(n-1)/2\rfloor$.
\end{theorem}
\begin{proof}
Here $n=2m+2$ with $m=t$ odd, since $\delta$ is odd.
Thus $m\geq5$. By Lemma~\ref{deficit} and the retained
Whiston theorem, the independent pair images have length
$r-2\leq m-1$. Failure of the desired rank bound would force
$r=m+1$. Whiston's exact equality clause then identifies
the pair image as the full $S_m$.

Write $\sigma=\rho_1$. The group $G_0=\gen{\sigma,K}$
is transitive on $B$, whereas $K$ fixes $b$, so
$\sigma=(b\ x)v$ for some $x\in X$ and an odd involution
$v$ supported on $X\setminus\{x\}$. Set $x'=x^\delta$
and $H=G_{>2}$. Its generators commute with $\sigma$ and
fix $b$, hence fix $x$; they also commute with $\delta$,
hence fix $x'$. Their $m-2$ pair images are independent,
as a subfamily of the images in Lemma~\ref{deficit},
and fix the pair $\{x,x'\}$. Whiston therefore makes their
image the full $S_{m-1}$ on the other pairs.

The only $H$-fixed points of $X$ are $x,x'$, because its
pair image moves every other pair. The permutation $v$
centralizes $H$ and fixes $x$, so also fixes $x'$.
On the remaining even number $m-1\geq4$ of pairs, the
accepted paired-row centralizer lemma makes $v$ even.
Its hypotheses require the actual full symmetric pair image
and the common pair swap, not transitivity of $H$ on points;
they allow both a transitive action and two faithful rows.
This contradicts the oddness of $v$ and proves the theorem.
\end{proof}

\paragraph{The remaining critical range.}
Combine this theorem with the separately written odd-tail
reduction, under its explicit smaller-degree alternating
induction. A counterexample in the singleton-endpoint branch
must have $n=2r\geq14$, $G_0=A_B$, and $K$ intransitive
on its support. Its residual $\delta$ has an odd number $t$
of transpositions and a \emph{nonempty even} number $f$ of
fixed points. Lemma~\ref{deficit} gives exactly
\[
 n=2+2t+f,\qquad r-2=t+f/2-1,\qquad
                         z\leq f/2+1.                 \tag{2}
\]
Since $K$ preserves both the nonempty fixed-point set and
the nonempty union of pairs, $z\geq2$. In particular $f=2$
forces $z=2$: $K$ is transitive on the two fixed points and
on the set of $t$ pairs, but need not be transitive on all
$2t$ points in those pairs. No IP is inferred for either
projected action.

\paragraph{Attribution and limit.}
The endpoint and pair-centralizer strategy is borrowed from
Cameron--Fernandes--Leemans--Mixer, \emph{Highest rank of a
polytope for $A_n$}, arXiv:1605.09173v1, Proposition~5.29,
physical page38, and its previously accepted full replacement
proof. The actual orbit-equivalence argument and inequality
(1) extend that proof to an intransitive tail without a
fixed-point-free hypothesis on the tail itself. Theorem~\ref{nofixed}
removes the tail-transitivity requirement while retaining the
necessary fixed-point-free residual hypothesis.
All retained Whiston inputs and the complete paid paired-row
proof remain dependencies. This is written mathematics; no
new computation, formalization or source-wide recovery is
claimed. The cases $f\geq2$, the source's further equality
graphs and the complete L044 recovery remain open here.
\end{document}
